%%%%%%%%%%%%%%%%%%%%%%%%%%%%%%%
%%%Autor: Briggette Román
%%%%%%%%%%%%%%%%%%%%%%%%%%%%%%%


\documentclass[12pt,a4paper]{book}
%%%%%%%%%%%%%%%%%
%Manejo de paquetes
%%%%%%%%%%%%%%%%%
\usepackage[utf8]{inputenc} 
\usepackage[spanish]{babel} 
\usepackage{hyperref}
\usepackage{algorithm}
\usepackage{algorithmic}
\floatname{algorithm}{Algoritmo}
\renewcommand{\listalgorithmname}{Lista de algoritmos}
\renewcommand{\algorithmicrequire}{\textbf{Entrada:}}
\renewcommand{\algorithmicensure}{\textbf{Salida:}}
\renewcommand{\algorithmicend}{\textbf{fin}}
\renewcommand{\algorithmicif}{\textbf{si}}
\renewcommand{\algorithmicthen}{\textbf{entonces}}
\renewcommand{\algorithmicelse}{\textbf{si no}}
\renewcommand{\algorithmicelsif}{\algorithmicelse,\ \algorithmicif}
\renewcommand{\algorithmicendif}{\algorithmicend\ \algorithmicif}
\renewcommand{\algorithmicfor}{\textbf{para}}
\renewcommand{\algorithmicforall}{\textbf{para todo}}
\renewcommand{\algorithmicdo}{\textbf{hacer}}
\renewcommand{\algorithmicendfor}{\algorithmicend\ \algorithmicfor}
\renewcommand{\algorithmicwhile}{\textbf{mientras}}
\renewcommand{\algorithmicendwhile}{\algorithmicend\ \algorithmicwhile}
\renewcommand{\algorithmicloop}{\textbf{repetir}}
\renewcommand{\algorithmicendloop}{\algorithmicend\ \algorithmicloop}
\renewcommand{\algorithmicrepeat}{\textbf{repetir}}
\renewcommand{\algorithmicuntil}{\textbf{hasta que}}
\renewcommand{\algorithmicprint}{\textbf{imprimir}} 
\renewcommand{\algorithmicreturn}{\textbf{devolver}} 
\renewcommand{\algorithmictrue}{\textbf{cierto }} 
\renewcommand{\algorithmicfalse}{\textbf{falso }} 
 % mi archivo de traducción
\usepackage{titling} %Para llamar los datos
\usepackage{titlesec} 
\usepackage{lipsum} %Muestras Lorem Ipsum
\usepackage{tikz}
\usepackage{afterpage}
\newcommand\blankpage{%
	\null
	\thispagestyle{empty}%
	\addtocounter{page}{0}%
	\newpage}
\usepackage{listings}
\renewcommand{\lstlistingname}{Código}
\renewcommand\lstlistlistingname{Índice de Código}

%%%%%%%%%%%%%%%%%%
%%%% Datos de la universidad, tesista y asesor
\title{{Esquema de Monitorización en Seguridad para entorno IoT }}
\author{Briggette Olenka Román Huaytalla}
\def\supervisor{Mg. Gipsy Miguel Arrunategui Angulo}
\def\university{\href{http://www.uni.edu.pe/}{Universidad Nacional de Ingeniería}}
\def\faculty{Facultad de Ciencias}
\def\school{Escuela Profesional de Ciencia de la Computación}

%%%%%%%%%%%%%%%%%%%
%%%%%%%%%%%%%%%%%%%
%%%%Inicio de archivo
\begin{document}

%%%%%%%%%%%%%%%%%%%
%%% Caratula

	\begin{titlepage}
		\begin{center}
		
		{\scshape \huge \university \par}
		\vspace{0.2cm}
		{\scshape \Large \faculty \par}
		\vspace{0.2cm}
		{\scshape \large \school \par}
%		\vspace{0.2cm}

		\begin{figure}[h]
			\centering
			\includegraphics[scale = 0.5]{figure/log_uni}
		\end{figure}
%		\vspace{0.2cm}
		
		{\LARGE \thetitle}\\[1cm]
		{\Large \textbf{SEMINARIO DE TESIS II}}\\[0.5cm]
		\large\emph{Autor: \\[0.3cm]}
		{\theauthor}\\[0.5cm]
		\large\emph{Asesor:\\[0.3cm] }
		{\supervisor }\\[1cm]
		{\large Diciembre, 2018}\\[4cm] 
		\end{center}
			
	\end{titlepage}
\afterpage{\blankpage}
%%%%%%%%%%%%%%%%%%

\frontmatter

%%%%%%%%%%%%%%%%%%
%%% Resumen
\titleformat{\chapter}[display]
{\Large\itshape}
{\filright\MakeUppercase{\chaptertitlename} }
{10ex}
{\titlerule\vspace{1ex}\filleft}
[\vspace{1ex}\titlerule]
\chapter{Resumen}
	
		\lipsum[1-2] 
	
	\vspace{10px}	
	\textbf{PALABRAS CLAVE:} IoT, Ciberseguridad, OWASP, vulnerabilidades, Gestión de Riesgos.
	
	\textit{El hardware es fácil de proteger: bloquearlo en una habitación, asegurarlo en un escritorio, o comprar uno de repuesto. La información posee más de un problema. Puede existir en más de un lugar, ser transportada a medio planeta en segundos, y ser robada sin que uno tenga conocimiento.}
	
	\begin{flushright}
		(Bruce Schneider, Protege tu Macintosh, 1994)
	\end{flushright}

%%%%%%%%%%%%%%%%%%%%%%%%%
%%% Indices y Listas

\afterpage{\blankpage}
\tableofcontents
\afterpage{\blankpage}
\listoffigures
\afterpage{\blankpage}
\listoftables
\afterpage{\blankpage}
\lstlistoflistings
\afterpage{\blankpage}
%%%%%%%%%%%%%%%%%%%%%%%%%%


%%%%%%%%%%%%%%%%%%%%%%%%%%
%%% Formato de capitulos
	
\titleformat{\chapter}[display]
{\bfseries\Large}
{\filright\MakeUppercase{\chaptertitlename} \Huge\thechapter}
{1ex}
{\titlerule\vspace{1ex}\filleft}
[\vspace{1ex}\titlerule]	
%%%%%%%%%%%%%%%%%%%%%%%%%%

\mainmatter
\chapter{Introducción}
\lipsum
\section{title}
\subsection{title}
\lipsum	
	
\end{document}